% --- Archon physics formalization source begin ---
% archon:physics
% archon:covers PhyXMiniProblems/problem_phyx_mini_0506.lean
% archon:source-report reports/phyx_mini/problem_phyx_mini_0506.source.json

\chapter{Physics problem phyx\_mini\_0506}
\label{ch:PhyXMiniProblems_problem_phyx_mini_0506}

\paragraph{Problem source.}
\#\# Physical scenario

Figure shows the wave function of a particle confined between \textbackslash{}( x = -4.0 \textbackslash{}text\{ mm\} \textbackslash{}) and \textbackslash{}( x = 4.0 \textbackslash{}text\{ mm\} \textbackslash{}). The wave function is zero outside this region.

\#\# Question

Calculate the probability of finding the particle in the interval \textbackslash{}( -2.0 \textbackslash{}text\{ mm\} \textbackslash{}leq x \textbackslash{}leq 2.0 \textbackslash{}text\{ mm\} \textbackslash{}).

\#\# Answer choices

- A: A: 0.19

- B: B: 0.27

- C: C: 0.13

- D: D: 0.36

\#\# Figure caption (auxiliary; use the image as primary evidence)

The image is a graph with labeled axes. 

- **Axes**:
  - The horizontal axis is labeled as \textbackslash{}(x \textbackslash{}, (\textbackslash{}text\{mm\})\textbackslash{}).
  - The vertical axis is labeled as \textbackslash{}(\textbackslash{}psi(x)\textbackslash{}).

- **Plot Details**:
  - A red line forms a piecewise linear function.
  - The line starts at the point (-4, -c), moves horizontally to (-4, c), then moves diagonally through the origin, and ends at the point (4, c).

- **Markers and Labels**:
  - Horizontal and vertical tick marks are present on both axes.
  - The constant value \textbackslash{}(c\textbackslash{}) is marked on the positive and negative sides of the vertical axis.

The graph represents a linear relationship with changes in slope at certain points on the \textbackslash{}(x\textbackslash{})-axis.

\#\# Dataset metadata

Quantum Phenomena; Numerical Reasoning, Physical Model Grounding Reasoning

\paragraph{Recorded answer/context.}
C: C: 0.13

\paragraph{Figure/image path.}
/root/proposal\_for\_physic/science-mango/phyx\_mini\_run/phyx\_data/test\_image/506.png

\paragraph{Formalization target.}
create a compiling Lean file with sorry bodies at `PhyXMiniProblems/problem\_phyx\_mini\_0506.lean`. The Lean declarations must preserve the physical quantities, dimensions or dimensional roles, figure labels, governing-law hypotheses, and final relation expressed by this problem.
Use Mathlib/Physlib names found through LeanExplore where available. If a physics API is missing, introduce faithful local abstractions rather than scalar placeholder aliases.

\begin{theorem}[Physics formalization target]
\label{thm:physics:phyx_mini_0506:target}
\lean{PhyXMini0506.central_interval_probability}
\uses{def:physics:phyx-mini-0506:phyxmini0506-leftwallmillimetres, def:physics:phyx-mini-0506:phyxmini0506-rightwallmillimetres, def:physics:phyx-mini-0506:phyxmini0506-querylowermillimetres, def:physics:phyx-mini-0506:phyxmini0506-queryuppermillimetres, def:physics:phyx-mini-0506:phyxmini0506-positionprobability, def:physics:phyx-mini-0506:phyxmini0506-roundstohundredth}
The assigned autoformalize agent should translate this physics problem into Lean declarations in the covered file, with theorem and lemma proof bodies written as `by sorry`.
\end{theorem}
\begin{proof}
This is an autoformalization task, not a proof task. Produce statements that can later be proved by the physics prover without weakening the physical model.
\end{proof}

% --- Archon declaration topology begin ---
\section{Declaration topology}
\begin{definition}[left Wall Millimetres]
  \label{def:physics:phyx-mini-0506:phyxmini0506-leftwallmillimetres}
  \lean{PhyXMini0506.leftWallMillimetres}
  The left wall of the confining region shown in the figure, in millimetres
\end{definition}
\begin{proof}
  This is the declared model definition of left Wall Millimetres.
\end{proof}
\begin{definition}[right Wall Millimetres]
  \label{def:physics:phyx-mini-0506:phyxmini0506-rightwallmillimetres}
  \lean{PhyXMini0506.rightWallMillimetres}
  The right wall of the confining region shown in the figure, in millimetres
\end{definition}
\begin{proof}
  This is the declared model definition of right Wall Millimetres.
\end{proof}
\begin{definition}[query Lower Millimetres]
  \label{def:physics:phyx-mini-0506:phyxmini0506-querylowermillimetres}
  \lean{PhyXMini0506.queryLowerMillimetres}
  The lower endpoint of the interval queried in the problem, in millimetres
\end{definition}
\begin{proof}
  This is the declared model definition of query Lower Millimetres.
\end{proof}
\begin{definition}[query Upper Millimetres]
  \label{def:physics:phyx-mini-0506:phyxmini0506-queryuppermillimetres}
  \lean{PhyXMini0506.queryUpperMillimetres}
  The upper endpoint of the interval queried in the problem, in millimetres
\end{definition}
\begin{proof}
  This is the declared model definition of query Upper Millimetres.
\end{proof}
\begin{definition}[position Probability]
  \label{def:physics:phyx-mini-0506:phyxmini0506-positionprobability}
  \lean{PhyXMini0506.positionProbability}
  The Born-rule probability of finding the particle between two position readouts. The integration variable is the numerical position in millimetres
\end{definition}
\begin{proof}
  This is the declared model definition of position Probability.
\end{proof}
\begin{definition}[Rounds To Hundredth]
  \label{def:physics:phyx-mini-0506:phyxmini0506-roundstohundredth}
  \lean{PhyXMini0506.RoundsToHundredth}
  RoundsToHundredth probability n means that the nearest-hundredth display of probability is n / 100, using the usual half-up convention at the lower boundary
\end{definition}
\begin{proof}
  This is the declared model definition of Rounds To Hundredth.
\end{proof}
% --- Archon declaration topology end ---
% --- Archon physics formalization source end ---
